% =====================================================================
\section{Introduction and Implementation Overview}
\label{sec:intro}
% =====================================================================

\subsection{Purpose of LipScribe}
LipScribe is an audio-visual speech recognition (AVSR) system that produces
text captions from a speaker's voice and lip movements together. Conventional
automatic speech recognition (ASR) performs well in quiet rooms but degrades
sharply when background noise, overlapping speakers or muffled speech corrupt
the audio signal. Lip movements are unaffected by acoustic noise, so a model
that also watches the mouth retains a usable source of evidence when the audio
becomes unreliable~\cite{li2023,afouras2022}. The implementation described here
therefore aims to build a recogniser whose accuracy degrades more gracefully in
noise than an audio-only system.

The proposal identifies deaf and hard-of-hearing people as the primary
beneficiaries, with lectures, meetings and public spaces as the intended
settings \spec. These users depend on captions being correct rather than merely
fast, so the implementation prioritises measurable transcription quality across
noise levels and treats latency as a constraint to be measured, not assumed.

\subsection{Problem statement and objectives}
The proposal states the central problem as: \emph{how can transcription be
improved by combining audio understanding with lip reading so that captions
remain usable even when the audio channel is severely degraded?} \spec. It
defines four objectives, each linked to a research question (RQ). \Cref{tab:objectives}
maps each objective to the part of the implementation that addresses it and
states whether that part belongs to the core system.

\begin{table}[H]
\centering
\caption{Mapping of proposal objectives to implementation components.}
\label{tab:objectives}
\small
\begin{tabularx}{\textwidth}{@{}p{3.4cm}X p{1.9cm}@{}}
\toprule
\textbf{Objective (RQ)} & \textbf{Implementation component} & \textbf{Scope} \\
\midrule
1. Disambiguate homophenes (RQ1) &
Cross-modal attention lets acoustic context resolve visually identical lip
shapes; temporal modelling over neighbouring frames (\cref{sec:model}). &
Core \\
2. Model temporal word sequences (RQ2) &
3D-CNN front-end, BiLSTM encoders and CTC training map frame sequences to
whole sentences (\cref{sec:visual,sec:audio,sec:model}). &
Core \\
3. Detect emotion from speaking style (RQ3) &
Separate landmark-based classifier trained on an emotion-labelled dataset
(\cref{sec:extensions}). &
Extension \\
4. Identify the spoken language (RQ4) &
Separate language-identification head trained on multilingual data
(\cref{sec:extensions}). &
Extension \\
\bottomrule
\end{tabularx}
\end{table}

\subsection{End-to-end workflow}
At run time the system executes the same sequence of stages for every segment
of input. \Cref{fig:workflow} summarises them; \cref{sec:architecture}
describes each stage as a software module.

\begin{enumerate}
  \item \textbf{Capture} synchronised video frames and audio samples from a
        camera and microphone, or read them from a recorded clip.
  \item \textbf{Detect the face} in every frame with MediaPipe Face
        Mesh~\cite{kartynnik2019,lugaresi2019} and locate the lip landmarks.
  \item \textbf{Extract the mouth sequence} by cropping a fixed-size region
        around the lips in each frame.
  \item \textbf{Extract audio features} (log-Mel spectrograms, optionally MFCCs)
        from the waveform.
  \item \textbf{Encode each modality} with its own neural encoder: a 3D-CNN and
        BiLSTM for video, and a BiLSTM for audio.
  \item \textbf{Fuse} the two encoded sequences with cross-modal attention.
  \item \textbf{Decode} the fused sequence into characters with a CTC output
        layer and beam search.
  \item \textbf{Display} the resulting text as a caption.
\end{enumerate}

\begin{figure}[H]
\centering
\begin{tikzpicture}[
  node distance=3mm,
  st/.style={draw=lsblue, rounded corners=2pt, fill=lsblue!6, align=center,
             font=\scriptsize\sffamily, minimum height=9mm, text width=13.5mm},
  ar/.style={-{Stealth[length=2mm]}, thick, lsblue}]
\node[st] (a) {Capture\\A/V};
\node[st, right=of a] (b) {Face\\mesh};
\node[st, right=of b] (c) {Mouth\\crop};
\node[st, right=of c] (d) {Audio\\features};
\node[st, right=of d] (e) {Encode\\streams};
\node[st, right=of e] (f) {Attention\\fusion};
\node[st, right=of f] (g) {CTC\\decode};
\node[st, right=of g] (h) {Caption\\display};
\foreach \x/\y in {a/b,b/c,c/d,d/e,e/f,f/g,g/h} \draw[ar] (\x) -- (\y);
\end{tikzpicture}
\caption{The eight stages of the LipScribe processing workflow.}
\label{fig:workflow}
\end{figure}

\subsection{Scope of the initial implementation}
The core deliverable is the audio-visual speech recogniser and a captioning
demonstration built on it. It covers objectives~1 and~2 and the evaluation
plan in the proposal: training on the GRID corpus~\cite{cooke2006} with
MUSAN noise~\cite{snyder2015} added at clean, 10, 5, 0 and $-5$\,dB SNR, and
reporting WER and CER on held-out speakers \spec.

Emotion detection and language identification (objectives~3 and~4) are not
part of the core architecture in this plan. GRID provides no emotion labels
and contains a single language, so neither task can be trained or evaluated
on it alone. \Cref{sec:extensions} describes how each could be investigated as
a separate module with its own data and evaluation, to be attempted only after
the core system works (\cref{sec:schedule}). Keeping them separate prevents
unvalidated extensions from delaying or obscuring the main result.

\subsection{Structure of this document}
\Cref{sec:architecture,sec:stack} define the system architecture and
technology stack. \Cref{sec:dataset,sec:visual,sec:audio} describe the data
and the two preprocessing pipelines. \Cref{sec:model,sec:training,sec:evaluation}
specify the model, its training and its evaluation. \Cref{sec:application}
covers the live captioning application, \cref{sec:extensions} the optional
extensions, and \cref{sec:schedule,sec:risks,sec:deliverables} the schedule,
risks and deliverables.
